\subsection{Nguồn dữ liệu và tổ chức PKLot}
PKLot được lựa chọn vì cung cấp ảnh bãi xe toàn cảnh cùng thông tin vùng ô và trạng thái chiếm dụng. Theo công bố của bộ dữ liệu, PKLot có 12.417 ảnh màu độ phân giải 1280 \(\times\) 720 pixel, thu nhận tại hai bãi xe ngoài trời trong các điều kiện nắng, nhiều mây và mưa. Nguồn dữ liệu và công trình giới thiệu PKLot được dẫn tại \cite{pklot,pklotdata}.

Archive sử dụng trong dự án có ba góc nhìn PUCPR, UFPR04 và UFPR05. UFPR04 và UFPR05 quan sát cùng bãi UFPR từ hai góc nhìn; PUCPR thuộc bãi khác. Ảnh được tổ chức theo camera, thời tiết và ngày. Dự án sử dụng ảnh toàn cảnh và XML để giữ bối cảnh cho detector; bộ ảnh từng ô đã cắt sẵn không thay thế đầu vào toàn cảnh.

\subsection{Nhãn ô đỗ và cấu hình hình học}
Mỗi phần tử space trong XML có mã ô, nhãn occupied và contour của vùng ô. Nhãn 1 biểu thị có xe, nhãn 0 biểu thị trống. Nhãn thiếu hoặc không hợp lệ không được suy đoán thành nhãn trống. Trạng thái chưa xác định là đầu ra mà hệ thống có thể sử dụng khi chưa đủ cơ sở quyết định; đó không phải lớp ground truth bổ sung của tập con đã chốt.

Tọa độ polygon được giữ trong ảnh gốc, gốc ở góc trên trái, trục x hướng sang phải và trục y hướng xuống dưới. Bộ JSON hình học của tập con có 5 layout: PUCPR có 1, UFPR04 có 3 và UFPR05 có 1. Mỗi frame liên kết đến layout bằng geometry\_id, còn nhãn trạng thái được lưu riêng theo mã ô. Cách biểu diễn này giữ được các biến thể chú thích thay vì ép mọi ảnh cùng camera dùng một polygon cố định.

\subsection{Kiểm tra chất lượng và bảo toàn metadata}
Với archive local đã kiểm tra bằng loader v2, dự án ghi nhận 12.416 XML và 1 ảnh thiếu XML. Trong số các frame có XML, 11.638 frame hợp lệ được giữ lại; 778 frame bị loại do nhãn hoặc mã ô không hợp lệ. Số lượng này là kết quả kiểm tra của dự án, khác với số ảnh được công bố cho toàn bộ PKLot. Danh sách lỗi được lưu để truy vết; không tự gán nhãn thay thế hoặc bỏ âm thầm ô lỗi trong một frame.

Loader hỗ trợ cả hai cách viết Point/point trong contour. Ngày thu nhận chuẩn được lấy từ timestamp hợp lệ trong tên file; ngày thư mục gốc được giữ trong directory\_date và nguồn ngày được ghi trong date\_source. Có 195 frame hợp lệ có ngày trong tên file khác ngày thư mục, bao gồm thư mục có tháng không hợp lệ. Dữ liệu gốc không bị đổi tên. Thời tiết Cloudy được chuẩn hóa thành overcast, đồng thời đường dẫn gốc vẫn được lưu.

Bảng~\ref{tab:valid-data} cho biết số frame hợp lệ trước khi chọn tập con thực nghiệm.

\begin{table}[htbp]
\centering
\caption{Số frame hợp lệ sau kiểm tra chất lượng}
\label{tab:valid-data}
\begin{tabular}{llr}
\toprule
Camera & Bãi xe & Frame hợp lệ \\
\midrule
PUCPR & PUCPR & 4.077 \\
UFPR04 & UFPR & 3.547 \\
UFPR05 & UFPR & 4.014 \\
Tổng &  & 11.638 \\
\bottomrule
\end{tabular}
\end{table}

\subsection{Tập con dùng chung và ngân sách xử lý}
Để phù hợp với ngân sách Colab và giữ so sánh công bằng, dự án cố định tập con pklot\_v1 dùng chung cho mọi phương pháp. Seed là 42; mỗi cặp camera/thời tiết có giới hạn tối đa 140 frame train, 30 validation và 30 test. Việc chọn frame luân phiên theo ngày và dùng thứ tự hash xác định; không lặp ảnh để đạt số lượng mong muốn.

Giao thức chính có tối đa 1.800 ảnh và thực tế chọn được 1.722 ảnh. Hai nhóm train UFPR04/rainy và UFPR05/rainy thiếu lần lượt 22 và 56 ảnh so với giới hạn; không lấy bù từ validation/test. Tập con ưu tiên khả năng so sánh giữa camera/thời tiết, nên phân bố này không đại diện tỉ lệ thời tiết tự nhiên của toàn bộ PKLot.

\subsection{Chia theo nhóm ngày và kiểm thử chuyển camera}
Train, validation và test được chia theo ngày thu nhận, không chia ngẫu nhiên từng frame hoặc từng crop ô. Một ngày chỉ thuộc một split trong mỗi giao thức, kể cả khi ngày đó có nhiều thời tiết hoặc camera; tất cả ô của một frame đi cùng split. Quy tắc này hạn chế việc ảnh gần nhau theo thời gian, chứa cùng xe đỗ, xuất hiện đồng thời trong tập học và tập kiểm thử.

Giao thức within\_camera đưa cả ba camera vào mỗi split và phân chia nhóm ngày xấp xỉ 70/15/15. Phương án chia được chọn trong 256 ứng viên xác định bởi seed, dựa vào số lượng và metadata camera/thời tiết, không dựa vào dự đoán hoặc điểm số test. Sau khi áp dụng giới hạn ảnh, số frame thực tế được trình bày trong Bảng~\ref{tab:splits}; tỉ lệ nhóm ngày không được hiểu là tỉ lệ ảnh chính xác.

Giao thức chuyển góc nhìn dùng UFPR04 cho train và validation và giữ UFPR05 cho test. Mọi ngày của camera đích được loại khỏi nguồn trước khi chia ngày nguồn xấp xỉ 80/20. Đây là kiểm thử góc nhìn mới trong cùng bãi UFPR, không phải kiểm thử chuyển sang một bãi xe độc lập. Hai giao thức là hai thí nghiệm riêng, không gộp manifest hoặc dùng model đã học camera đích để báo kết quả chuyển camera.

\begin{table}[htbp]
\centering
\caption{Các tập con thực nghiệm đã khóa của dự án}
\label{tab:splits}
\begin{tabular}{lrrr}
\toprule
Giao thức & Train & Validation & Test \\
\midrule
Trong camera & 1.182 & 270 & 270 \\
Chuyển UFPR04 \(\rightarrow\) UFPR05 & 420 & 90 & 90 \\
\bottomrule
\end{tabular}
\end{table}

\subsection{Vai trò của các split và khả năng tái lập}
Train được sử dụng cho các thành phần cần học, như HOG/SVM. Validation được dùng để chọn tham số, padding, overlap và ngưỡng suy trạng thái. Với detector pretrained, checkpoint và cấu hình được chốt trước khi đánh giá test. Sau khi khóa lựa chọn, tất cả phương pháp chạy trên cùng tập test của giao thức tương ứng; không chọn lại tham số bằng kết quả test.

Các manifest train/val/test lưu đường dẫn tương đối, camera, bãi xe, thời tiết, ngày, mã layout và nhãn từng ô. File split\_lock.json ghi cấu hình, danh sách ngày và SHA256 các artifact; notebook tải cùng bộ metadata từ GitHub thay vì tự chia lại trên mỗi máy. Bộ kiểm tra xác nhận không giao ngày giữa các split và chỉ có camera đích trong test của thí nghiệm chuyển camera. Bộ pklot\_v1 là giao thức của dự án, không phải bản sao split chính thức của PKLot.
